\section{Background}

\subsection{Last Value Prediction}

The simplest guess for an instruction's result is the value it
produced the last time it ran. Last Value Prediction (LVP) uses
exactly this guess. It is a computational
predictor that applies the identity function to the last observed
value~\cite{ravenel2025architecting}. Predictions are usually
associated with static instructions using their PCs, so observations
from earlier executions can be reused.

Instruction values may change across executions, so LVP uses a
confidence counter to decide when a prediction is admitted. The design
in \cite{ravenel2025architecting} uses a two-bit saturating counter,
incremented on correct predictions and decremented on incorrect
ones. A prediction is admitted only when the counter is saturated.
In this implementation, the admission threshold is configurable.
The counter is updated for every eligible load that completes,
whether or not a prediction was made or used. This means entries
below the threshold can keep gaining confidence.

Accuracy alone is not enough, and neither is coverage. A predictor with high accuracy but few
predictions may give little benefit. A predictor with high coverage
may not help if the predicted instructions are not on critical
dependency chains.

LVP is especially relevant to in-order processors, where a
long-latency load can delay younger dependent instructions. A
predicted value lets eligible consumers, including dependent memory
operations, proceed before the producer completes. Memory operations
that would otherwise wait for one another can then overlap.

The producer still runs, and its real value is checked against the
prediction when it completes. If an
incorrect value was used, the processor must recover the affected
younger instructions to keep the architectural state correct.

\subsection{ASTRA Address Prediction}

Besides values, the paper also predicts load addresses, using a
predictor called ASTRA (Address Stride Alias)~\cite{ravenel2025architecting}. It uses an order-1 differential
finite-context method (D-FCM), with two modifications: it uses the
last observed stride directly as history, and its tables are untagged.

ASTRA is indexed by load PC. Its Last Address Table (LAT) stores
the last committed virtual address and signed stride for each
entry. The stride indexes a second table, Stride$\rightarrow$Stride
(SS), which gives the stride that is added to the LAT address to
form the predicted address. When a load retires, ASTRA computes the
new stride, updates the corresponding SS entry, and records the
committed address and stride in the LAT.

Because the SS table is untagged, loads that see the same stride
share entries. The paper uses this aliasing on purpose. When one
load updates an entry, other loads can use it to predict a new
access pattern. As a result, ASTRA favors prediction
coverage over accuracy, and it depends on address validation before
a prediction is used (Section~\ref{sec:bg-vastra}).

The original ASTRA predicts the complete load address. V-ASTRA
instead predicts the load's base-register value and adds the
instruction displacement. This project implements the base-address
form, so its LAT stores the base-register value rather than the
complete effective address (Section~\ref{sec:impl-astra}). ASTRA
predicts an address, not the loaded data. Getting the load value
requires V-ASTRA's separate value cache.

\subsection{V-ASTRA and the Value Cache}
\label{sec:bg-vastra}

To supply the data itself, V-ASTRA (Value-Address Stride Alias)
uses ASTRA's predicted address to read a dedicated value cache, which
stores previously observed load values. In the paper, the value cache holds 4\,KiB of data
in 16-byte entries. Loads allocate entries at commit, while stores
update matching entries without allocating new ones.

V-ASTRA predicts a load's base address and adds the displacement
from the decoded instruction to get its effective address.
The paper checks the predicted base at issue by comparing it
with the actual base-register value. A wrong prediction is
therefore rejected before it is used, so V-ASTRA needs neither a
confidence counter nor recovery from value mispredictions.

For a bypass to be safe, the value cache must also contain the
requested bytes, and no older in-flight store may modify them. The
paper uses a counting Bloom filter to detect stores that might
conflict, and an invalidation FIFO to handle older stores that
issue after a prediction. When all checks succeed, the load skips
the L1D access and its dependents can issue in the same cycle~\cite{ravenel2025architecting}.

The implementation splits the check into two stages and replaces the Bloom
filter and FIFO with an exact scan of older stores
(Section~\ref{sec:impl-vcache}).

\subsection{Dependency Handling}

When an instruction needs a value from an earlier instruction, it
has to wait for that producer. In an in-order processor,
a consumer waiting for a long-latency load can also delay younger
independent instructions, because issue follows program order. Under
stall-on-use, a dependent instruction cannot issue until the
producer's value becomes available.

Two mechanisms can reduce this delay. Dependency release
makes the consumer eligible to issue earlier but does not provide
the required value. Value forwarding supplies the producer's
value directly, so an eligible consumer can proceed before the
producer completes. In the paper, issue-time V-ASTRA validation
lets dependents issue in the same cycle as the load. MinorCPU
handles operands differently. Ordinary non-memory instructions read
their operands at commit, while memory operations can issue to
the LSQ earlier (Section~\ref{sec:bg-minor}). For this reason, the
forwarding mechanism in this implementation is restricted to eligible
dependent loads.

\subsection{META Selection}

Some loads have predictable addresses, while others have predictable
values. ASTRA/V-ASTRA targets the first kind and LVP the second, so
META combines them to improve
prediction coverage \cite{ravenel2025architecting}.

In the original design, both predictors operate in parallel, and
META selects a prediction at issue time. A validated V-ASTRA value
is preferred when the predicted base is correct and the requested
data is available in the value cache. Otherwise, META uses an
LVP prediction if its confidence counter is saturated. If neither
is available, the load is not predicted. Both predictors are updated
when the load retires~\cite{ravenel2025architecting}.

V-ASTRA
validates its address before use, whereas LVP remains speculative
until the actual value is available. So an incorrect LVP prediction
needs recovery only if it was used. The paper replays younger
instructions through the scoreboard, while the implementation squashes and
refetches them (Section~\ref{sec:impl-meta}).

\subsection{gem5 MinorCPU}
\label{sec:bg-minor}

All of this work is done in MinorCPU, gem5's in-order processor
model~\cite{gem5}. Its four pipeline stages are Fetch1 (instruction fetch),
Fetch2 (instruction separation and branch prediction), Decode, and
Execute. Execute contains the functional-unit pipelines, in-flight
instruction queue, commit logic, and load-store queue (LSQ).

Instructions issue to functional units in program order when their
source registers are ready according to the scoreboard. The
scoreboard tracks expected result availability and the latest
in-flight writer of each register, but does not hold register values.
Results are written to architectural registers at commit, also in
program order.

Memory instructions can compute their addresses and enter the LSQ
before they reach the head of the in-flight queue. However, they
issue one at a time, and ordinarily only after the instructions that
produce their sources have committed. Loads write their destination
registers at commit. Branches are resolved in Execute, and a
misprediction causes younger wrong-path instructions to be discarded.

The memory functional unit uses \texttt{extraAssumedLat} to estimate
when dependents may issue. With the default setting, dependents may
issue after an assumed load latency, but they still wait for the
producer to commit before executing. Setting this latency to zero
enables the stall-on-use configuration, in which dependents wait for
commit before issuing.
